\section{Introducción}\label{sec:intro}
El Reto AgroCebada 2026, organizado por FIRA, pide estimar el rendimiento de grano de cebada de 59 parcelas, es decir, cuántas toneladas de grano se cosecharon
por cada hectárea sembrada, a partir
de imágenes satelitales, clima y topografía. Para entrenar se cuenta con otras 138 parcelas del mismo ciclo cuyo
rendimiento sí se conoce. En todo el documento el rendimiento se expresa en toneladas por hectárea (t/ha). Por
ejemplo, una parcela de 5 hectáreas que rinde 4~t/ha produjo 20 toneladas de grano.

El objetivo de este documento es que todo el equipo tenga el mismo contexto antes de construir el modelo. Explicamos
qué pide el reto, cómo se produce la cebada en la región, qué datos se entregaron y qué tipo de problema estadístico
es. Después describimos lo que encontramos en una primera revisión de los datos y planteamos las hipótesis que guían
el trabajo.

\textbf{Tres tipos de afirmaciones.} Para que quede claro qué tan firme es cada cosa, distinguimos entre
\emph{hechos}, que vienen de las bases del reto o se leen directamente en los archivos, \emph{observaciones
preliminares}, que salen de cálculos rápidos y todavía deben comprobarse, e \emph{hipótesis}, que son ideas por
probar con el modelo.
